\documentclass[12pt,a4paper]{article}
\usepackage[utf8]{inputenc}
\usepackage[spanish,es-noquoting]{babel}
\usepackage[T1]{fontenc}
\usepackage{lmodern}
\usepackage{geometry}
\geometry{top=2cm, bottom=2cm, left=2.5cm, right=2.5cm}
\usepackage{xcolor}
\usepackage{listings}
\usepackage{enumitem}
\usepackage{fancyhdr}
\usepackage{hyperref}
\usepackage{tikz}
\usetikzlibrary{arrows.meta, positioning, calc, shapes.geometric}

\definecolor{codegreen}{rgb}{0,0.6,0}
\definecolor{backcolour}{rgb}{0.95,0.95,0.92}

\lstset{
    language=C,
    backgroundcolor=\color{backcolour},
    commentstyle=\color{codegreen},
    keywordstyle=\color{blue},
    basicstyle=\ttfamily\small,
    breaklines=true,
    frame=single,
    numbers=left,
    tabsize=4,
    extendedchars=true,
    showstringspaces=false,
    inputencoding=utf8,
    literate={á}{{\'a}}1 {é}{{\'e}}1 {í}{{\'i}}1 {ó}{{\'o}}1 {ú}{{\'u}}1 {ñ}{{\~n}}1 {¡}{{\textexclamdown}}1 {¿}{{\textquestiondown}}1
}

\pagestyle{fancy}
\fancyhf{}
\lhead{Monitorías Sistemas Operativos 2026-I}
\rhead{Capítulo 4: Variables de Condición}
\cfoot{\thepage}

\title{\textbf{Guía de Aprendizaje: Variables de Condición}}
\author{Malak Sanchez \\ \small Monitorías de Sistemas Operativos}
\date{\today}

\begin{document}

\maketitle

\section{Concepto}
Las \textbf{Variables de Condición} permiten que un hilo se suspenda (se duerma) hasta que una condición específica se cumpla. A diferencia de un Mutex, que protege datos, las variables de condición permiten que los hilos se comuniquen sobre el \textit{estado} de esos datos.

\section{Regla de Oro}
Una variable de condición \textbf{siempre} debe usarse en conjunto con un Mutex para evitar condiciones de carrera entre la comprobación de la condición y el acto de dormirse.

\begin{center}
\begin{tikzpicture}[font=\small, >=Stealth,
    node/.style={draw, thick, minimum width=2.5cm, minimum height=1cm, align=center}]

    \node[node, fill=orange!20] (t1) at (-4, 0) {Hilo Consumidor \\ (Espera señal)};
    \node[node, fill=green!20] (cond) at (0, 0) {Variable de \\ Condición};
    \node[node, fill=blue!20] (t2) at (4, 0) {Hilo Productor \\ (Envía señal)};

    \draw[->, ultra thick] (t1) -- (cond) node[midway, below] {wait()};
    \draw[->, ultra thick] (t2) -- (cond) node[midway, below] {signal()};
    
    \node at (0, -1.5) {\textit{``Avísame cuando el tanque esté lleno''}};

\end{tikzpicture}
\end{center}

\newpage

\section{Funciones Principales}
\begin{lstlisting}[numbers=none]
#include <pthread.h>

pthread_cond_t cond;
pthread_mutex_t lock;

// Esperar (libera el mutex y se duerme)
pthread_cond_wait(&cond, &lock);

// Despertar a UN hilo
pthread_cond_signal(&cond);

// Despertar a TODOS los hilos
pthread_cond_broadcast(&cond);
\end{lstlisting}

\section{Ejemplo: Productor-Consumidor (Básico)}
\begin{lstlisting}
#include <stdio.h>
#include <pthread.h>

int combustible = 0;
pthread_mutex_t lock;
pthread_cond_t cond;

void* reabastecer(void* arg) {
    pthread_mutex_lock(&lock);
    combustible += 40;
    printf("Reabastecido. Combustible: %d\n", combustible);
    pthread_cond_signal(&cond);
    pthread_mutex_unlock(&lock);
    return NULL;
}

void* conducir(void* arg) {
    pthread_mutex_lock(&lock);
    while (combustible < 40) {
        printf("No hay suficiente. Esperando...\n");
        pthread_cond_wait(&cond, &lock);
    }
    combustible -= 40;
    printf("Viaje realizado. Restante: %d\n", combustible);
    pthread_mutex_unlock(&lock);
    return NULL;
}

int main() {
    pthread_t t1, t2;
    pthread_mutex_init(&lock, NULL);
    pthread_cond_init(&cond, NULL);

    pthread_create(&t1, NULL, conducir, NULL);
    pthread_create(&t2, NULL, reabastecer, NULL);

    pthread_join(t1, NULL);
    pthread_join(t2, NULL);

    pthread_mutex_destroy(&lock);
    pthread_cond_destroy(&cond);
}
\end{lstlisting}

\section{Ejercicios}
\begin{enumerate}
    \item ¿Por qué se usa un \texttt{while} en lugar de un \texttt{if} al llamar a \texttt{pthread\_cond\_wait}?
    \item Implemente una cola sincronizada (Buffer Circular) usando variables de condición para manejar los estados de ``Cola Llena'' y ``Cola Vacía''.
\end{enumerate}

\end{document}
